\label{sec:conclusoes}

Os codificadores de passagem de mensagens sobre o grafo dos canais ativos, com os escalares do evento, reconstroem o vértice de depósitos de \qty{10}{\MeV} (média de \qty{530.1}{\pe} por evento de teste) com erro 3-D médio de \num{25.95} a \qty{26.27}{\cm} e RMSE combinado de \num{18.55} a \qty{18.80}{\cm} nos quatro melhores. Com capacidade equiparada, o erro 3-D médio dos \num{11} codificadores varia em até \qty{2.89}{\cm}, e esses quatro não se distinguem nos \num{937} eventos de teste. O GATv2-cascata, escolhido pelo menor erro 3-D médio de teste, tem RMSE por eixo, no referencial dos sensores, de \num{20.45}, \num{14.35} e \qty{20.46}{\cm} em $x$, $y$ e $z$.

Nos mesmos eventos e na luz limpa, o GATv2-cascata tem erro 3-D médio \qty{1.71}{\cm} menor que o do XGBoost e RMSE combinado \qty{2.30}{\cm} menor; é melhor nos eventos escuros e na cauda, mas pior no quartil mais claro. Em $z$, o baricentro calibrado tem RMSE \num{2.07} vezes o do GATv2-cascata por causa dos vértices perto da tampa, onde o erro do baricentro é sistemático, viés que o APEX não encontrou na sua simulação. Abaixo de $z = {}$\qty{20}{\m}, o baricentro fica próximo do GATv2-cascata (\num{23.29} e \qty{22.54}{\cm}), e seu erro é sobretudo estatístico: com todos os fótons, cai para \qty{8.33}{\cm}. O~RMSE combinado do GATv2-cascata é \qty{10.6}{\percent} menor que o da rede da etapa anterior, calculado a partir dos valores publicados por eixo, de outra partição e da geometria anterior às correções; a diferença está em $x$ e $y$, pois em $z$ o valor daquela rede, \qty{20.87}{\cm}, fica dentro do intervalo de \qty{95}{\percent} do GATv2-cascata. Em luz detectada igual com o rendimento médio, esse RMSE combinado é \qty{51.0}{\percent} menor que a largura combinada do APEX (\qty{41.3}{\percent} com o mínimo e \qty{28.0}{\percent} em energia igual), média quadrática de larguras sem definição publicada, lidas de um gráfico de outra simulação. Em~$z$, contra \qty{20.46}{\cm} do GATv2-cascata, o APEX tem \qty{24.8}{\cm} com o rendimento médio, acima do intervalo de \qty{95}{\percent}, e \qty{19.4}{\cm} com o mínimo, dentro dele; só em energia igual, em que detecta \num{3.4} vezes mais luz, é melhor (\qty{13.7}{\cm}).

{\clubpenalty=10000 O erro do GATv2-cascata concentra-se nos eventos escuros, longe das paredes longas (\qty{42.79}{\cm} de erro 3-D médio no quartil mais escuro, contra \qty{11.67}{\cm} no mais claro), mas não se separou o efeito da estatística de fótons do da posição. Os resultados valem para uma energia, uma eficiência, uma geometria, um sorteio de fótons por evento e vértices na grade do treino. Cada codificador foi treinado uma vez, sem a busca de hiperparâmetros feita nos \textit{boosters}, com lotes e passos diferentes; não houve ablação dos \num{33} escalares; os \textit{boosters} foram treinados em luz com aumento; e os eventos de teste já tinham orientado o desenho final dos codificadores.\par}

{\clubpenalty=10000 Como trabalhos futuros, propõe-se treinar e avaliar os codificadores em várias eficiências, energias e geometrias; gerar vários sorteios de fótons por evento e vértices fora da grade; repetir a comparação com orçamento de otimização fixo, busca de hiperparâmetros por codificador, \textit{boosters} ajustados e treinados na luz limpa e várias sementes, com e sem os escalares, em eventos de teste novos; prever as coordenadas a partir dos estados dos canais, para testar se resumir o grafo num só vetor limita a precisão; combinar \textit{boosters} e codificadores, que erram menos em faixas de luz diferentes; comparar com o APEX numa simulação equivalente, confirmando a correspondência dos eixos e a definição da largura; e medir o tempo de inferência por evento.\par}
